% Vervang de TODO door de bevestigde methode van de betreffende medestudent.
\subsection{RQ1: De foto's vooraf bewerken}
Als het programma geen hand vindt, proberen we het nog eens met een bewerkte versie van dezelfde foto. De medestudent onderzoekt verschillende bewerkingen. Daarna moet worden nagekeken of de gevonden punten werkelijk bij de hand passen. De voorgestelde bewerkingen staan in bijlage A.
\drafttodo{PREP}{Vul in welke bewerkingen worden getest en hoe wordt gecontroleerd of de gevonden handpunten bruikbaar zijn.}

\subsection{RQ2: Met meer foto's oefenen}
De tweede medestudent laat het programma met steeds meer verschillende foto's oefenen. De toetsfoto's blijven hetzelfde. Anders kan een hogere score ook komen doordat de toets gemakkelijker is geworden. Een bewerkte kopie telt niet als een nieuwe oorspronkelijke foto.
\drafttodo{SIZE}{Vul in met hoeveel foto's wordt geoefend en hoe vaak de vergelijking wordt herhaald.}
